\section{混合重正化方案}
\label{sec:hybrid}

如上一节所述，LaMET 展开从计算大动量 $P^z$ 下、整个距离范围
$-\infty<z<\infty$ 内的坐标空间关联函数 $\tilde h(z,P^z)$ 出发。
在格距为 $a$ 的离散格点上，定义 \eq{esme} 中 $\tilde h(z,P^z)$ 的非定域夸克双线型算符可以相乘性地重正化为~\cite{Ji:2017oey,Ishikawa:2017faj,Green:2017xeu}
\begin{align}\label{eq:ren}
&\left[\bar{\psi}(z)\Gamma W(z,0)\psi(0)\right]_{\rm B} \nn\\
&= e^{\delta m|z|}Z(a) \left[\bar{\psi}(z)\Gamma W(z,0)\psi(0)\right]_{\rm R}\,,
\end{align}
（除格点效应外）~\cite{Constantinou:2017sej,Chen:2017mie}。此处左边的算符由裸场与裸耦合定义，以下标``B''标记；右边的算符已重正化，以下标``R''标记。除非明确声明，我们总是假定重正化后的关联在作到动量空间的傅里叶变换之前最终定义于 $\MS$ 方案。这里既有与 $z$ 无关的对数发散，也有随 $z$ 变化的线性发散：前者源于夸克与胶子场以及 Wilson 线端点顶点的重正化，包含在因子 $Z(a)$ 中；后者来自 Wilson 线自能，被吸收进指数因子 $e^{\delta m|z|}$，$\delta m$ 即``质量修正''。

文献中已有许多方案~\cite{Ishikawa:2016znu,Chen:2016fxx,Monahan:2016bvm,Radyushkin:2017cyf,Constantinou:2017sej,Green:2017xeu,Stewart:2017tvs,Braun:2018brg,Li:2020xml} 用于重正化上述格点关联函数 $\tilde h(z, a, P^z)$，其中 RI/MOM 方案最为常用~\cite{Constantinou:2017sej,Stewart:2017tvs}。在该方法中，人们先在固定规范、深欧氏态（动量平方 $-p^2\gg\Lambda_{\rm QCD}^2$）中计算双定域算符 $O(z,a)$ 的矩阵元（截腿格林函数）$Z(z,-p^2,a)$，然后定义 $\overline{\rm MS}$ 算符为
\begin{align}
             O_{\overline{\rm MS}}(z,\mu) \equiv Z_{\overline{\rm MS}}(z,-p^2,\mu)\frac{O(z,a)}{Z(z,-p^2,a)},
\end{align}
其中 $Z_{\overline{\rm MS}}$ 把 RI/MOM 重正化的结果转换到 $\overline{\rm MS}$ 方案。两个 $Z$ 因子的规范依赖与 $-p^2$ 依赖相互抵消；右边具有无发散的正常连续极限 $a\to 0$。

然而，尽管 RI/MOM 方法对定域算符是合理的，用于非定域算符时却有潜在问题。例如当 $z$ 大时，$Z_{\overline{\rm MS}}(z,-p^2,\mu^2)$ 含有 $z$ 的 IR 对数，其 $z$ 依赖的微扰计算并不可靠。此外，虽然 RI/MOM 因子 $Z(z,-p^2,a)$ 有助于抵消格点 UV 发散，但大 $z$ 处的复合算符同样包含非微扰物理。因此两个 $Z$ 因子都含有无法相消的非微扰效应，会改变 $O(z)$ 的 IR 性质。所以 RI/MOM 重正化方案在大 $z$ 处并不可靠。再者，涉及胶子分布时它需要外部离壳胶子态，从而带来与规范变异算符潜在混合的复杂性~\cite{Wang:2019tgg}。

除大距离的重正化问题外，短距离的重正化也有微妙之处。对双定域算符作标准重正化可使其在任何非零 $z$ 处有限，但在 $z\to 0$ 极限下变为发散。另一方面，若对小 $z$ 处的大对数作重求和，结果在 $z=0$ 处为零。然而格点在 $z=0$ 的结果趋于矢量流的矩阵元。显然，$a\to 0$ 与 $z\to 0$ 这两个极限不可交换。

为解决这些问题，本节提出一种重正化关联函数的混合方案。该方案的要点是：把短距离与长距离关联分开并分别重正化，再在中间距离 $z_{\rm S}$ 处将两套程序匹配起来。匹配点必须位于 $[0,z_{\rm LT}]$ 之内，即关联算符的领头扭度（LT）近似成立的区域。关于 $z_{\rm LT}$ 取值的讨论见第 V 节。

\subsection{短距离 $0\le |z|\le z_{\rm S}$ 的重正化}
\label{sec:short}

要重正化 $0\le |z|\le z_{\rm S}$ 的 $\tilde h(z,a, P^z)$，须特别关注关联函数在 $z\to 0$ 极限下的行为。

在连续 $\MS$ 方案中，$z\rightarrow 0$ 极限并不光滑，会出现额外的对数型 UV 发散 $\sim \ln z^2$，重求和后其结果为零。然而格点矩阵元 $\tilde h(z,a, P^z)$ 的情形并非如此。对有限格距与非零 $z$，$\tilde h(z,a, P^z)$ 含有与裸场波函数重正化相关的 UV 发散，形如
$\alpha_s(a)\ln (z^2/a^2)$。
在小 $z$ 处（尤其当 $z=0$ 或 $a$ 时），$\tilde h(z,a, P^z)$ 存在离散化效应，并与格点正规化的定域矩阵元 ${\bar\psi}\Gamma \psi$ 相联系。特别是当 $\Gamma=\gamma^\mu$ 时，${\bar\psi}\gamma^\mu \psi$ 守恒，其矩阵元在 $a\to 0$ 极限下有限。展示格点调节子与小物理距离之间这一有趣相互作用的函数是 $\ln [(z^2+a^2)/a^2]$。

小 $z$ 区域的上述分歧可以通过格点正规化与连续 $\MS$ 方案之间的微扰转换来消除，但这已知收敛缓慢。更有效的策略是通过格点重正化消去 $\ln z^2$ 依赖，这对应于不同于 $\MS$ 的方案``X''。只要 $z$ 位于领头扭度区 $|z|\le z_{\rm S}$（$z_{\rm S}$ 小于 $z_{\rm LT}$），X 方案与 $\MS$ 之差便可用微扰论计算。

例如，X 方案可以通过构造 $\tilde h(z,a, P^z)$ 与同一算符 $O(z,a)$ 的另一矩阵元之比来实现，
\begin{align}\label{eq:ratio}
 \frac{\tilde h(z,a,P^z)}{Z_X(z,a)}\,,~~~~~~{\rm for\ }|z|\le z_{S}\,,
\end{align}
其中重正化因子 $Z_X$ 对应不同的矩阵元选择。$Z_X$ 的可能选择包括：
\begin{itemize}
\item{$O(z,a)$ 在单粒子深欧氏态中、固定于某特定规范（如 Landau 规范）下的截腿格林函数，它定义了 RI/MOM 型方案~\cite{Constantinou:2017sej,Stewart:2017tvs,Alexandrou:2017huk,Chen:2017mzz}。}
\item{$O(z,a)$ 在 $P^z=0$ 强子态中的矩阵元，视应用而定~\cite{Radyushkin:2017cyf,Orginos:2017kos}。}
\item{$O(z,a)$ 的真空（$|\Omega\rangle$）期望值~\cite{Braun:2018brg,Li:2020xml}。}
\end{itemize}
第二与第三种选择的矩阵元是规范不变的，因而无需固定规范。第三种选择要求算符的量子数与真空相同。如上所述，$z_{\rm S}$ 必须小于 $z_{\rm LT}$，后者在第 V 节估计约为
$0.25 \sim 0.33$~fm。当然，应显式验证最终结果对 $z_{\rm S}$ 小变化的稳定性。

由于算符 $O(z,a)$ 具有相乘可重正化性，\eq{ratio} 的比值中所有 UV 发散相消，从而可以取连续极限，
\begin{align}\label{eq:ratio2}
\lim_{a\to0} \frac{{\tilde h}(z,a,P^z)}{Z_X(z,a)} = \frac{{\tilde h}(z,\epsilon,P^z)}{Z_X(z,\epsilon)} \equiv\tilde h^X(z,P^z)\,,
\end{align}
其中第一个等号之后的项指同一比值的 $\MS$ 计算，$\epsilon$ 对应连续理论中维数正规化 $d=4-2\epsilon$。在 $|z|\to0$ 极限下，$\ln z^2$ 依赖与外部态无关，故在比值中相消，使后者在 $z=0$ 处有限。此外，在领头扭度区 $z\le z_{\rm S} \le z_{\rm LT}$，我们可以通过坐标空间因子化公式把任意 X 方案的比值微扰地匹配到 LF 关联 $h(\lambda,\mu)$~\cite{Izubuchi:2018srq,Radyushkin:2017lvu}
\begin{align}\label{eq:coordfact}
\tilde h^X(\lambda, P^z) & =\int_{0}^1 d\alpha\ {\cal C}^X\Big(\alpha,\frac{\lambda^2\mu^2}{(P^z)^2}\Big) h(\alpha\lambda,\mu) \nn\\
&\qquad + {\cal O}(z^2\Lambda_{\rm QCD}^2)\ ,
\end{align}
其中 ${\cal C}^X$ 是匹配系数，其在 X 方案（如 RI/MOM）中对重正化标度的依赖已被隐去；高扭度贡献在上式中亦被压低。

第二种 $Z_X$ 选择的一圈匹配系数已在文献 ~\cite{Zhang:2018ggy,Radyushkin:2017lvu,Izubuchi:2018srq} 中得到，RI/MOM 方案的系数则可由 ~\cite{Constantinou:2017sej,Izubuchi:2018srq} 提取。第二与第三种选择的两圈结果已作为级数展开在 ~\cite{Li:2020xml} 中计算。

\eq{coordfact} 在动量空间还有等价形式~\cite{Izubuchi:2018srq}，匹配系数已在一圈~\cite{Stewart:2017tvs,Liu:2018uuj,Liu:2018hxv} 与两圈~\cite{Chen:2020ody} 阶计算得到。
由 $P^z=0$ 矩阵元定义的 $Z_X$ 的两圈匹配系数可从 ~\cite{Chen:2020ody,Li:2020xml} 提取。

\subsection{大距离 $z\ge z_{\rm S}$ 的重正化}
\label{sec:long}

在大距离 $z\ge z_{\rm S}$ 处，UV 重正化需要仔细评估，因为 RI/MOM 与比值方案都会引入不合意的非微扰效应。唯一不会引入此类额外非微扰物理的重正化途径，是显式地分别减除线性发散（或 $\delta m$）与对数发散~\cite{Chen:2016fxx,Ji:2017oey,Ishikawa:2017faj,Green:2017xeu}；原则上这可以借助辅助场方法实现~\cite{Green:2017xeu,Green:2020xco}。

要计算 Wilson 线的质量重正化 $\delta m$，文献中有许多建议。这里我们给出一份未必完备的清单：

\begin{itemize}
\item{可以在大 $z$ 处拟合强子矩阵元，其主导衰变为
\begin{equation}\label{eq:mass1}
  \tilde h(z,a, P^z) \sim \exp(-\delta m |z|) \ .
\end{equation}
$\delta m$ 可以通过在大 $z$ 处把比值 $\ln (\tilde h(z+a,a, P^z)/\tilde h(z,a, P^z))$
拟合为 $z$ 的常数来获得。当然，结果必须与 $P^z$ 无关，例如可取 $P^z=0$。或者，也可以在所有 $z$ 上把 $\ln \tilde h(z,a, P^z)$ 对 $1/a$
依赖进行拟合。该方法尚未用真实格点数据加以研究。}
\item{可以像在 RI/MOM 重正化因子中那样使用单夸克格林函数
\begin{align}
          Z(z,-p^2,a)&=\int d^4x d^4y\ e^{ip\cdot (x-y)}\nn\\
          &\times  \langle \Omega|T\psi(x) O(z,a) \bar{\psi}(y)|\Omega\rangle\,,
\end{align}
其渐近行为为 $Z(z) \sim \exp(-\delta m |z|)$。此矩阵元需要固定规范，
已在 ~\cite{Izubuchi:2019lyk,Huo:2019vdl} 中研究过。}
\item{也可以使用 $O(z,a)$ 的真空矩阵元
\begin{equation}
        S(z) = \langle \Omega|O(z,a)|\Omega\rangle
\end{equation}
它是规范不变的。$S(z)$ 在很大 $z$ 处同样有行为
$S(z) \sim \exp(-\delta m |z|)$。~\cite{Braun:2018brg,Li:2020xml} 考虑过这一做法。
}
\item{规范不变的 Polyakov 圈也给出两重夸克之间的静态势。关于该途径的文献很多，参见例如 ~\cite{Appelquist:1977tw,Schroder:1999xf,Smirnov:2009fh,Philipsen:2002az,Jahn:2004qr}。}
\item{还可以直接在固定规范中计算 Wilson 线 $W(z)$ 的真空期望值，大距离处同样有 $\langle W(z) \rangle \sim \exp(-\delta m |z|)$。~\cite{Green:2017xeu,Green:2020xco} 用辅助场方法考虑过这一做法。}
\end{itemize}
值得指出的是，尽管上述建议原则上都可行，它们在格点实现时会遇到不同的实际问题，因而有些可能比其他的更好用。

质量重正化 $\delta m$ 是规范无关的，正如夸克的极点质量一样~\cite{Kronfeld:1998di}。对上述无需固定规范的建议而言，这显然成立。在需要固定规范的情形下，可以通过构造适当的规范不变算符来证明任何其他规范中的结果都相同~\cite{Philipsen:2001ip}。尽管规范无关，$\delta m$ 会依赖于蒙特卡罗模拟所用的具体作用量以及上述矩阵元的定义。在真空矩阵元的情形下，$\delta m$ 在某些规范中可以被诠释为非微扰极点质量~\cite{Philipsen:2001ip}。

由上述各矩阵元算得的 $\delta m$ 对格距 $a$ 有如下依赖关系
\begin{equation}
   \delta m = m_{-1}(a)/a + m_0 \ ,
\end{equation}
其中 $m_{-1}(a)$ 是幂次发散的系数，与具体矩阵元无关。与 $a$ 无关的项 $m_0$
来源较为复杂，可能来自以下几个方面：
\begin{itemize}
\item{Renormalon 效应：原则上 $m_{-1}(a)$ 可以微扰计算，领头阶为 $2\pi\alpha_s(a)/3$，类似于 Wilson 费米子表述中的质量反项。然而微扰级数并不收敛：当截断于 $n\sim 1/\alpha_s(a)$ 阶时，级数具有 $a\Lambda_{\rm QCD}$ 量级的不确定度，从而给 $m_0$ 带来贡献~\cite{Ji:1995tm,Beneke:1998ui,Bauer:2011ws,Bali:2013pla}。
这意味着在非微扰拟合中，$m_{-1}(a)$ 只能确定到 $a\Lambda_{\rm QCD}$ 量级的不确定度。因此预期存在 $\Lambda_{\rm QCD}$ 量级的额外 $m_0$ 贡献。}
\item{极点质量：对某些矩阵元（如双定域算符的真空期望值），其 $z$
依赖可以视为源于由一个无穷重夸克与一个轻夸克组成的介子的极点质量。此时 $\delta m$ 就是除去线性发散后的极点质量。}
\item{有限 $P^z$ 效应：有限 $P^z$ 的关联函数具有长程关联
$\exp(-\lambda/\xi(P^z))$，其中 $\xi(P^z)$ 是关联长度。这一贡献计入 $m_0$。}
\item{拟合效应：由于数据总处于有限的 $z$ 与 $a$，指数衰减未必总能与代数衰减区分开，因此还存在计入 $m_0$ 以及影响 $m_0$ 与 $m_{-1}$ 分离的拟合不确定度。}
\end{itemize}
总之，$m_0$ 依赖于所用的格点矩阵元与拟合流程~\cite{Zhang:2017bzy,Chen:2017gck,Izubuchi:2019lyk,Huo:2019vdl}。

作为例子，图~\ref{fig:deltam} 给出了由夸克 RI/MOM 重正化因子确定的 $m_{-1}$ 与 $m_0$ 的数值：重正化因子在标度 $\mu_R=1.8~ {\rm GeV}$、$p^z_R=0$ 下计算，使用了 MILC 合作组四个系综（$a\approx \{0.045, 0.06, 0.09, 0.12\} {\rm fm}$、π 介子质量 $310$ MeV）~\cite{Bazavov:2012xda}。受第~\ref{sec:regge} 节将研究的大 $z$ 渐近行为的启发，我们用如下简化形式
\begin{equation}
e^{-(\frac{m_{-1}}{a}+m_0)|z|}\frac{c_1}{|z|^{d_1}}
\end{equation}
来拟合四种不同格距下的重正化因子。值得指出，$m_{-1}$ 从 $O(\alpha_s)$ 开始，因此我们在拟合中也计入了耦合对 $a$ 的依赖。对 $a\approx 0.12$ fm，系数 $m_{-1}$ 与 $m_0$ 的拟合结果为
\begin{equation}
m_{-1}=0.234\pm 0.012, \ \ m_0=(350\pm 60)\, {\rm MeV}.
\end{equation}

\begin{figure}[tbp]
\centering
\includegraphics[width=0.9\columnwidth]{images/npr_deltam_new}
\vspace*{-2em}
\caption{对四个系综（格距 $a\approx \{0.045, 0.06, 0.09, 0.12\} {\rm fm}$、π 介子质量 $310$ MeV）计算的夸克 RI/MOM 重正化因子的拟合，组态来自 MILC 合作组~\cite{Bazavov:2012xda}。}
\label{fig:deltam}
\end{figure}

原则上可以选择只减除幂次发散部分，即 $m_{-1}/a$，把确定程度较差的 $m_0$ 项留在格点矩阵元中，动量展开会处理其余部分。事实上，减除不同 $m_0$ 所带来的差异是 $O(1/P^z)$ 效应，~\cite{Xiong:2013bka} 已对此作了微扰证明。更确切地说，假设有两个准 LF 关联，它们定义准 PDF 且彼此相差一个因子 $e^{-m|z|}$（$m \sim \Lambda_{\rm QCD}$）：
\begin{equation}
\tilde h_1(z,P^z)=\tilde h_2(z,P^z)e^{-m|z|} \ ,
\end{equation}
那么傅里叶变换到动量空间后二者满足
\begin{align}
&\tilde f_1(y,P^z)-\tilde f_2(y,P^z)\nonumber \\
&=\frac{1}{\pi}\int_{-\infty}^{\infty} dy' \frac{\delta}{\delta^2+(y-y')^2} \left[\tilde f_2(y',P^z)-\tilde f_2(y,P^z) \right]\ ,
\end{align}
其中 $\delta=m/P^z$。若诸 $\tilde f$ 平方可积且一阶导数连续，则可证明当 $\delta \rightarrow 0$ 时
\begin{align}
|\tilde f_1(y,P^z)-\tilde f_2(y,P^z)|\sim \delta \ .
\end{align}
因此不同方案的模糊性在大 $P^z$ 极限下消失。

然而这仍不能令人满意：由于线性发散的非微扰效应，LaMET 展开似乎将成为 $M/P^z$ 幂次的展开而非 $(M/P^z)^2$ 的展开，这会显著降低收敛速度。这里我们考虑克服这一缺陷的可能途径。回顾一下：$(M/P^z)^2$ 的 LaMET 展开是在不存在线性发散的 $\MS$ 方案中进行的；在一般方案中该展开或许含有 $1/P^z$ 的奇数次幂。因此，可以选择 $m_0$ 使得 $\overline{\rm MS}$ 方案的条件
\begin{equation}
      \delta m_{\overline{\rm MS}} = 0
\end{equation}
由非微扰计算得到满足。
我们把这样的减除质量取值记为
\begin{equation}
\delta m_c = m_{-1}/a + m_{0c}\ .
\end{equation}
其中 $m_{0c}$ 可如下确定：在小 $z$（约 $z_S$）处且 QCD 微扰论适用的区域，把格点上的矩阵元 $\tilde h_1(z, P^z, a)$ 与 $\overline{\rm MS}$ 结果匹配。另一策略是在 $\Lambda_{\rm QCD}$ 附近的某个范围内变动 $m_{0}$，找出使 $\tilde f(y, P^z)$ 中 $1/P^z$ 线性项消失的值 $m_{0c}$。
这类似于寻找 Wilson 费米子出现类似幂次发散裸夸克质量的临界值 $\kappa_c$~\cite{Lepage:1992xa}。

质量减除后的算符 $O(z,a)e^{-\delta m(a)|z|}$ 不含幂次发散，但对 $a$ 仍保留对数依赖。这些残余的对数发散与 $z$ 无关，原则上可以用辅助场方法重正化~\cite{Green:2017xeu,Green:2020xco}。不过实践中更方便的做法是：直接令短距离区域与长距离区域在 $z=z_{\rm S}$ 处的 $O(z,a)$ 重正化矩阵元相互匹配，以此固定重正化常数 $Z_{\rm hybrid}(a)$——这本质上是一个连续性条件，
\begin{align}
Z_{\rm hybrid} e^{-\delta m |z_{\rm S}|}\langle P| O(z_{\rm S},a) | P \rangle  = \frac{\langle P| O(z_{\rm S},a) | P \rangle }{Z_X(z_{\rm S},a) }\,,
\end{align}
由此得
\begin{align}
Z_{\rm hybrid}(z_{\rm S},a) = e^{\delta m |z_{\rm S}|}/{Z_X(z_{\rm S},a)  }\,.
\end{align}
这样一来只需计算 $\delta m$。当然，还需变动 $z_{\rm S}$ 以检验最终结果是否稳定。

长距离区域的匹配系数 ${\cal C}_{\rm hybrid}$ 与 X 方案的匹配系数相关联。例如，若采用 $P^z=0$ 矩阵元作重正化~\cite{Zhang:2018ggy,Radyushkin:2017lvu,Izubuchi:2018srq}，则
\begin{align}\label{eq:hybridm1}
&{\cal C}_{\rm hybrid}(\alpha, z^2\mu^2, z^2/z_{\rm S}^2) = {\cal C}_{\rm ratio}(\alpha, z^2\mu^2) \nn\\
&\qquad + \delta(1-\alpha){\alpha_sC_F\over 2\pi}{3\over 2} \ln {z^2\over z_{\rm S}^2}\ \theta(|z|-z_{\rm S})\,.
\end{align}
然而，由于 $z^2\mu^2$ 与 $z^2/z_{\rm S}^2$ 的对数，上述匹配系数仅对 $|z| \ll \Lambda_{\rm QCD}^{-1}$ 有效，否则就必须将 $\alpha_s$ 演化到高度非微扰的区域，以重求和 $|z|\sim \Lambda_{\rm QCD}^{-1}$ 处的大对数。鉴于我们的最终目标是对结果作傅里叶变换以获得 PDF，这会引入不可控制的系统误差。

为了以更清晰的方式划分微扰与非微扰区域，可以在动量空间进行匹配：只要 $yP^z$ 足够大，就没有任何东西妨碍使用大 $z$ 关联。原则上应先把混合方案转换为 $\MS$ 方案——因子化公式正是在该方案中被证明的~\cite{Ma:2014jla,Ji:2020ect}——即在坐标空间使用转换因子
\begin{align}\label{eq:hybrid3}
    {\cal Z}(z,z_S,\mu) & = {Z_{\MS}\over Z_{\rm hybrid}} =\frac{\theta(z_{\rm S}\!-\!|z|)}{Z^{\overline{\rm MS}}_{X}(z,\mu)}
    +\frac{\theta(|z|\!-\! z_{\rm S})}{Z^{\overline{\rm MS}}_{X}(z_S,\mu)}\,,
\end{align}
其中 $1/Z^{\overline{\rm MS}}_{X}$ 把``$X$''方案转换到 $\MS$，而对 \eq{hybridm1} 有
\begin{align}\label{eq:zmsr}
    Z^{\overline{\rm MS}}_{\rm ratio} = 1 - {\alpha_sC_F\over 2\pi}\left[{3\over 2}\ln{z^2\mu^2\over 4e^{-2\gamma_E}}+ {5\over2}\right]\,.
\end{align}
转换因子 ${\cal Z}$ 对所有 $z$ 都是微扰的，因为它在 $|z|>z_{\rm S}$ 的大距离不含 $\ln(z^2)$。随后便可对 $\MS$ 准 LF 关联作傅里叶变换，并在动量空间把它匹配到 PDF。

既然方案转换对所有 $z$ 都是微扰的，我们也可以先做傅里叶变换，直接把混合方案的准 PDF 匹配到 PDF；匹配系数 $C_{\rm hybrid}$ 由 \eq{hybridm1} 的双重傅里叶变换给出：
\begin{align}\label{eq:hybridm2}
&C_{\rm hybrid}(\xi, \mu^2/(p^z)^2, z_{\rm S}^2\mu^2) = C_{\rm ratio}(\xi, \mu^2/(p^z)^2) \nn\\
&\quad+{\alpha_sC_F\over 2\pi}{3\over 2} \left[-{1\over |1-\xi|_+}+ \frac{2 \text{Si}((1-\xi) \lambda_{\rm S})}{\pi  (1-\xi)} \right]\,,
\end{align}
其中 $C_{\rm ratio}$ 见 ~\cite{Izubuchi:2018srq}，$\xi=y/x$，$\lambda_{\rm S}=z_{\rm S}p^z$ 而 $p^z=xP^z$ 为部分子动量。plus 函数定义为
\begin{align}
	{1\over |1-\xi|_+} &\equiv \lim_{\beta\to0^+}\left[ {\theta(|1-\xi|-\beta)\over |1-\xi|} + 2\delta(1-\xi)\ln\beta\right] \,.
\end{align}

我们还可以推导出短距离重正化采用 RI/MOM 方案时相应的方案转换因子与匹配系数。在 $-z^2_{\rm S}p^2 \ll 1$ 的极限下，$|z|<z_{\rm S}$ 的 RI/MOM 重正化因子 $Z(z,-p^2)$ 等于零动量矩阵元所对应的因子，故结果与 \eqs{hybrid3}{hybridm2} 相同。但若 $-z^2_{\rm S}p^2$ 有限，则需要利用 ~\cite{Constantinou:2017sej,Stewart:2017tvs} 的结果来推导方案转换因子与匹配系数。
最后，既然 PDF 结果必须与格点重正化方案无关，我们可以尝试不同的短距离方案并检验它们彼此的一致性。

在动量空间中，匹配系数含有 $\mu/(yP^z)$ 的对数，当 $y\sim \Lambda_{\rm QCD}/P^z$ 时它变为非微扰的。这与幂计数参量 $\Lambda^2_{\rm QCD}/(yP^z)^2$ 一致。因此系统不确定性的性质是清楚的：只有推向更高的 $P^z$ 才能改进小 $x$ 处的精度。
